\documentclass[11pt]{article}
\usepackage[margin=1in]{geometry}
\usepackage{amsmath,amssymb,amsthm,mathtools}
\usepackage[hidelinks]{hyperref}
\allowdisplaybreaks[2]
\setlength{\emergencystretch}{3em}
\newtheorem{theorem}{Theorem}[section]
\newtheorem{lemma}[theorem]{Lemma}
\newtheorem{proposition}[theorem]{Proposition}
\theoremstyle{remark}\newtheorem{remark}[theorem]{Remark}
\newcommand{\E}{\mathbb E}
\newcommand{\dd}{\,d}
\newcommand{\one}{\mathbf1}
\newcommand{\Gp}{G_+}
\newcommand{\Cl}{\mathcal C_q}
\title{A pointwise response-ledger cap and residual closure\\S0, Route 3}
\author{Analytical increment for review}
\date{Version 1; October 6, 2026}
\begin{document}\maketitle
\begin{abstract}
The response ledger bounds the sign-good amplitude at every time
inside the residual bootstrap. This supplies the missing cap at an
interior first-exit endpoint and closes RCET's Volterra argument.
No independent backward clock estimate or terminal amplitude cap is
assumed. The result is conditional on the existing original-entry,
strong-learning and secondary-gain contracts, and on explicitly
listed instantaneous core, chart and outside budgets. A bounded-chart
mass estimate derives the required small chart mass from the transverse
entry budget. No Stage 2 transfer or further passage is undertaken.
\end{abstract}

\section{Setting, partition, and retained inputs}
Every new label has prefix \texttt{an03lcrcv1:}.
Use the exact balanced original flow, with
\[
 P_1=N(e_1,q^2I_2),\quad P_2=N(\rho e_2,q^2I_2),\quad
 P=(P_1+P_2)/2,\quad \rho\in[13/20,7/10],\quad
 \lambda=\rho^2,\quad Q=\lambda+q^2.
\]
The equations and energies are those of CLOCK, BULK and RCET:
\[
 U'=R_\theta^TW,\quad W'=R_\theta U,\quad
 R_\theta=\E_P[(X-f)X^T\one_{\{a_\theta^TX>0\}}],\quad
 |U|^2=|W|^2=m.
\]
For any fixed cohort $S$, set
\begin{align*}
 z&=(U_2+W_2)/2, & a&=(U_2-W_2)/2,\\
 H_S&=\int_Sz^2\dd\lambda_0, & A_S&=\int_Sa^2\dd\lambda_0,\\
 E_S&=H_S+A_S=\tfrac12\int_S(U_2^2+W_2^2)\dd\lambda_0,
 &J_S&=\tfrac12\int_S(U_1^2+W_1^2)\dd\lambda_0.
\end{align*}
Balance gives $\mu_S:=\int_Sm=E_S+J_S$.
Write $c_S=Q^{-1}\E_{P_2}[X_2(f_S)_2]$, $c=c_{\rm all}$ and
$I(t)=Q\int_{t_0}^t(1-c)$.

Take a fixed, disjoint partition into core $\mathcal K$, sign-good
group $\Gp$, bounded-chart group $\mathcal D$, secondary group
$\Sigma$, and outside group $\mathcal B$.
The union defining $\Gp$ is counted once: it contains the clock
$\Cl$ and RCET's primary $\eta=0$ labels.
Set $G=\Gp\cup\mathcal D\cup\Sigma$, $H=H_{\Cl}$, and
\[
 Y=\sqrt{J_G}/q,\qquad
 r(t)=\sup_\theta|R_{11}(t,\theta)|,\qquad
 \ell(t)=\sup_\theta\max(|R_{12}|,|R_{21}|).
\]
No group is a moving ``currently charted'' mask. Every label is
charged once. The theorem retains the following inputs on
$[t_0,t_b]$; none is a residual-norm or clock-defect assumption.

\paragraph{Instantaneous budgets and the prescribed interval.}
The response satisfies $0\le c\le c_b<1$ and
$t_b-t_0\le L_T\log(1/q)$. The strong core satisfies
\begin{equation}\label{an03lcrcv1:core}
 |\theta|+|\psi|\le Rq\quad\hbox{on }\mathcal K,
 \qquad \mu_{\mathcal K}\le M_{\mathcal K},\qquad
 \sup_{[t_0,t_b]}\mu_{\mathcal B}/q\longrightarrow0.
\end{equation}
The fixed chart group satisfies
\begin{equation}\label{an03lcrcv1:chart}
 |\theta|+|\psi|\le Zq\quad\hbox{on }\mathcal D
 \quad(t_0\le t\le t_b),\qquad Z\ge1.
\end{equation}
This is a retained chart-state budget, not a claim of chart retention
from an entry cap. A displacement cap about a bounded reference
is included by increasing $Z$ to include that reference.

\paragraph{Existing S1 contracts.}
The original joint entry satisfies
\begin{equation}\label{an03lcrcv1:entry}
 J_G(t_0)\le C_Jq^3,\qquad E_{\Gp}(t_0)=o(1),\qquad
 H(t_0)\asymp q^{k},\quad k=10(1-\lambda).
\end{equation}
Uniformly on $\Gp$ at entry, $z>0$, $|a|/z\le1/4$ and
$qj/z\le C_0$, where $j=((U_1^2+W_1^2)/2)^{1/2}$.
On $\Cl$ the stronger bound $qj/z\le C_{\rm clk}q^{1/8}$ holds.
These are original-flow entry contracts, not target-only inferences.
The fixed constant $C_J$ comes from the entire explicit cohort,
so it is independent of its later subdivision at $C_0$.
With $\kappa_1=\int_{\mathcal K}W_1U_1\dd\lambda_0$, retain S1's
\begin{equation}\label{an03lcrcv1:learning}
 \int_{t_0}^{t_b}|1-\kappa_1|\dd t\le B_s.
\end{equation}
Initialization and the foundational dissipation identity give a
uniform bound on $\mathcal L(0)$, hence on all diagonal residual
norms; it is not an additional hypothesis.

\paragraph{Existing S1/S2 secondary contracts.}
Retain BULK Lemma 5.1 and Proposition 5.2's original-entry radial
profile, bounded reference, pre-exit distance envelope, and actual
mass-gain bound at every intermediate endpoint. Explicitly, with
$\dd\eta=m(t_0,\alpha)\dd\lambda_0$, entry displacement
$\xi_\alpha\le Z_\Sigma$, and $h_1=q^{10d-2s}$,
\begin{align}
 \eta([0,Z_\Sigma])&\le Cq^3,\quad
 \eta\{\xi>u\}\le Cq^3\min\{1,(h_1/u)^{p_2}\},
 \label{an03lcrcv1:profile}\\
 D_\alpha(t)&\le C_D\xi_\alpha e^{\beta I(t)/2}
 \quad\hbox{until the first }qD_\alpha=\delta_0,
 \label{an03lcrcv1:distance}\\
 m(t,\alpha)&\le C_m m(t_0,\alpha)
 \max\{1,C_gq^2\xi_\alpha^2e^{\beta I(t)}\}^{1+\nu}.
 \label{an03lcrcv1:gain}
\end{align}
Here $0<\beta\le1$, $0\le\nu\le1/100$,
$d=(1-\lambda P_\rho)/2$, $s=2d/(1-\lambda)$ and
$p_2=(1-3\lambda/2)/d$, as in BULK. The reference has both strong
coordinates bounded by a fixed $R_\Sigma$; $D_\alpha$ is the joint
scaled displacement from that same-field reference.
The same contracts must apply to the chosen fixed secondary group,
with no omitted labels or moving-mask flux. A fixed entry cap alone
does not supply \eqref{an03lcrcv1:distance} or \eqref{an03lcrcv1:gain}.

\section{Response ledger, sign convention, and absolute charges}
\begin{lemma}[Exact ledger and one-sided negative-gate estimate]
\label{an03lcrcv1:ledger}
Define the nonnegative negative part $v_-:=(-v)_+$ and
\[
 \tau_S=Q^{-1}\int_S W_2\E_{P_2}[X_2(U\cdot X)_-]\dd\lambda_0.
\]
For every fixed cohort,
\begin{equation}\label{an03lcrcv1:identity}
 c_S=H_S-A_S+\tau_S.
\end{equation}
If $W_2\ge0$ on $S$, without any input-angle restriction,
\begin{equation}\label{an03lcrcv1:taulower}
 \tau_S\ge-C\mu_S\exp[-\rho^2/(4q^2)].
\end{equation}
For arbitrary signs,
\begin{equation}\label{an03lcrcv1:responsecharge}
 |c_S|\le E_S+\frac{2q}{\sqrt Q}\sqrt{E_SJ_S}.
\end{equation}
\end{lemma}
\begin{proof}
The identity $v_+=v+v_-$ and the diagonal second moment of $P_2$
give $\E_{P_2}[X_2(U\cdot X)_+]=QU_2+
\E_{P_2}[X_2(U\cdot X)_-]$. Integrate against $W_2/Q$ and use
$U_2W_2=z^2-a^2$. This is exactly CLOCK's convention
$\tau_S=Q^{-1}\E[X_2\int_S W_2(-U\cdot X)_+]$;
rearrangement gives its $H_S=c-c_{S^c}-\tau_S+A_S$.

When $W_2\ge0$, only $X_2<0$ can make the tail term negative.
Since $|W_2||U|\le m$, Cauchy--Schwarz gives
\begin{align*}
 \tau_S&\ge-\frac{\mu_S}{Q}
 \E_{P_2}[|X_2||X|\one_{\{X_2<0\}}]\\
 &\ge-\frac{\mu_S}{Q}
 (\E_{P_2}[X_2^2|X|^2])^{1/2}
 \mathbb P(X_2<0)^{1/2}
 \ge-C\mu_S e^{-\rho^2/(4q^2)}.
\end{align*}
The fourth moment is uniformly bounded, $Q\ge169/400$, and the
Gaussian Chernoff bound is $\mathbb P(X_2<0)\le e^{-\rho^2/(2q^2)}$.
No receiver gate or projected angular margin occurs in this proof.

For \eqref{an03lcrcv1:responsecharge}, the positive part is
1-Lipschitz, so replacing $(U_1X_1+U_2X_2)_+$ by $(U_2X_2)_+$
costs at most $|U_1X_1|$. Moreover
$|\E X_2(U_2X_2)_+|\le Q|U_2|$ and
$\E|X_1X_2|\le q\sqrt Q$. Consequently
\[
 |c_S|\le\int_S|W_2U_2|+\frac q{\sqrt Q}\int_S|W_2U_1|
 \le E_S+\frac{2q}{\sqrt Q}\sqrt{E_SJ_S}.
\]
All integrands are dominated by a constant times $m|X|^2$;
finite mass justifies the label/expectation interchange.
\end{proof}

\begin{lemma}[Pointwise cap]
\label{an03lcrcv1:pointcap}
Suppose $z>0$, $|a|\le z/3$ on $\Gp$ at the time considered.
Writing ${\rm rest}=\mathcal D\cup\Sigma$, one has
\begin{equation}\label{an03lcrcv1:cap}
 \frac89 H_{\Gp}\le c_b+E_{\rm rest}
 +\frac{2q}{\sqrt Q}\sqrt{E_{\rm rest}J_{\rm rest}}
 +Cq^2+C\mu_{\mathcal B}
 +C\mu_{\Gp}e^{-\rho^2/(4q^2)}.
\end{equation}
\end{lemma}
\begin{proof}
Here $W_2=z-a\ge2z/3>0$ and $A_{\Gp}\le H_{\Gp}/9$.
Apply \eqref{an03lcrcv1:identity}--\eqref{an03lcrcv1:taulower}
and $c=c_{\Gp}+c_{\rm rest}+c_{\mathcal K}+c_{\mathcal B}\le c_b$.
Use \eqref{an03lcrcv1:responsecharge} on the combined rest group.
The core has $E_{\mathcal K}\le C_Rq^2\mu_{\mathcal K}$ and
$J_{\mathcal K}\le\mu_{\mathcal K}$, so its absolute response is
$O(q^2)$. Finally $|f_{\mathcal B}(X)|\le\mu_{\mathcal B}|X|$
gives $|c_{\mathcal B}|\le C\mu_{\mathcal B}$.
For a different fixed bound $|a|\le y_{\max}z$, $y_{\max}<1$,
the coefficient $8/9$ becomes $1-y_{\max}^2$.
\end{proof}

\section{Verification on stopped intervals}
\begin{lemma}[Local applicability and derived clock estimates]
\label{an03lcrcv1:localclock}
Fix $B,M$. On any $[t_0,t_*]\subset[t_0,t_b]$ where the cross
residual sum is at most $Mq$ and $\int_{t_0}^{t_*}r\le B$,
BULK Theorem 4.1 and RCET Theorem 2.1 apply with that endpoint.
For sufficiently small $C_0=C_0(B,M)>0$ and then small $q$,
\begin{align}
 z&>0,\quad |a|/z\le1/3\quad\hbox{on }\Gp,
 \qquad E_{\Gp}\le\tfrac{10}{9}H_{\Gp},
 \label{an03lcrcv1:signgood}\\
 \int_s^t|\varepsilon_{\Gp}|&\le D_+\le1,
 \qquad \int_s^t|\varepsilon_{\Cl}|
 \le D_{\rm clk}=O(q^{1/8}+q^2\log(1/q))=o(1),
 \label{an03lcrcv1:defects}\\
 H_{\Gp}(s)&\le e H_{\Gp}(t)e^{-b_0(t-s)},
 \qquad b_0=Q(1-c_b),\quad t_0\le s\le t\le t_*.
 \label{an03lcrcv1:backward}
\end{align}
\end{lemma}
\begin{proof}
The entry conditions and fixed-label definitions are anchored at
$t_0$. Response and residual norms are pointwise conditions on the
interval being used. The nonnegative integral $\int r$ and the
duration bound restrict to every initial subinterval. Diagonal norms
come from $\E_{P_i}|e|^2\le4\mathcal L(0)$. None of these source
theorems requires information after $t_*$, an endpoint amplitude
cap, or a pre-existing clock defect. Their constants are independent
of the chosen endpoint.

Use the conservative comparison
$y\le1/4+2LX/(3b_*)$ in choosing RCET's smaller admissible $C_0$;
it still gives $y\le1/3$. RCET provides the uniform individual
defect on all of $\Gp$ after making
$C_{\rm clk}q^{1/8}\le C_0$. Its weighted average has
$D_+\le C(B,M)(C_0+q^2\log(1/q))\le1$.
BULK with $\eta=1/8$ separately gives the sharper clock estimate.
Integrating the exact logarithmic identity on $[s,t]$ proves
\eqref{an03lcrcv1:backward}. Positivity and domination are exactly
those of the reviewed source theorems, so no new division at zero
or varying-label operation is introduced.
\end{proof}

\begin{lemma}[Small mass of the fixed chart group]
\label{an03lcrcv1:chartmass}
On the stopped interval of Lemma \ref{an03lcrcv1:localclock}, the
retained chart budget and transverse entry budget imply
\begin{equation}\label{an03lcrcv1:chartmassbound}
 \mu_{\mathcal D}(t)\le4C_Je^{2B}q^3,\qquad
 E_{\mathcal D}(t)\le C Z^2e^{2B}q^5
\end{equation}
for sufficiently small $q$, after $Z,M,B$ have been fixed.
In particular $Z\mu_{\mathcal D}=o(1)$ and
$Z^2q^2\mu_{\mathcal D}=o(1)$.
\end{lemma}
\begin{proof}
Loss gives a fixed bound $|R_{22}|\le C_{\mathcal L}$ and the
residual hypothesis gives $\ell\le C Mq$ by BULK's entrywise
residual estimate. For $U=\sqrt m\,u$, $W=\sqrt m\,w$,
\[
 (\log m)'=2w^TR_\theta u
 \le2r+Cq^2(MZ+Z^2),
\]
because $|u_2|,|w_2|\le Zq$ on $\mathcal D$.
Thus $\mu_{\mathcal D}(t)\le
 e^{2B+Cq^2(MZ+Z^2)L_T\log(1/q)}\mu_{\mathcal D}(t_0)$.
For small $q$, both first direction components have squared size
at least $1/2$, so $\mu_{\mathcal D}(t_0)\le2J_G(t_0)$.
Make the extra exponential at most two to obtain
\eqref{an03lcrcv1:chartmassbound}. The energy estimate follows from
$E_{\mathcal D}\le Z^2q^2\mu_{\mathcal D}$.
This derives the mass bound; only continued bounded-chart geometry
remains an input. No unproved strong-stage mass law is used.
\end{proof}

\begin{lemma}[Secondary energy without a future cap]
\label{an03lcrcv1:secondary}
On the same stopped interval, also suppose $E_{\Gp}<E_g$ for a
fixed $E_g$. Under \eqref{an03lcrcv1:profile}--\eqref{an03lcrcv1:gain},
\begin{equation}\label{an03lcrcv1:secondarybound}
 E_\Sigma(t)\le C_\Sigma(q^5+q^{\Gamma_\nu}H(t)),\qquad
 \Gamma_\nu\ge\frac{289169}{924000}>\frac3{10}.
\end{equation}
The constant can depend on $E_g$ and the fixed S1/S2 and geometric
constants, but not on $q$ or $t_*$.
\end{lemma}
\begin{proof}
Check the hypotheses of BULK's secondary implication on this very
interval. The original entry profile and seed are retained S1
contracts. The bounded reference, first-exit distance envelope and
all-endpoint gain bound are retained S2 contracts. The response
guard supplies $I'\ge0$ and $I(t_0)=0$. Lemma
\ref{an03lcrcv1:localclock} supplies the clock defect using only the
stopped residual bound. There is no need to import a terminal cap:
before its last simplification, BULK proves
\[
 E_\Sigma\le C\{q^5+q^\gamma H+q^{\Gamma_\nu}H^{1+\nu}\},
 \quad \gamma=1-5\lambda+\frac{2\lambda}{1-\lambda},
 \quad \Gamma_\nu=\gamma-\nu[10(1-\lambda)-2].
\]
Here $e^{I(t)}\le e^{D_{\rm clk}}H(t)/H(t_0)$ is available on
$[t_0,t]$. The algebraic inequality
$H\le H_{\Gp}\le E_{\Gp}<E_g$ holds even before proving
$y\le1/3$. Hence $H^{1+\nu}\le E_g^\nu H$ and
$q^\gamma\le q^{\Gamma_\nu}$. This proves
\eqref{an03lcrcv1:secondarybound} using BULK's already proved rational
margin. At entry its profile and bounded coordinates directly give
$E_\Sigma(t_0)=O_{Z_\Sigma,R_\Sigma}(q^5)$, without any clock input.
\end{proof}

\section{Joint closure and its constant order}
\begin{theorem}[Route 3; conditional original-flow residual closure]
\label{an03lcrcv1:closure}
Under the retained inputs of Section 1, choose the constants in
the order specified below. For sufficiently small $q$, on the whole
prescribed passage interval,
\begin{align}
 \|e_1\|_{P_1}+\|e_2\|_{P_2}&\le C,
 &\|e_2\|_{P_1}+\|e_1\|_{P_2}&\le Cq,
 \label{an03lcrcv1:residualconclusion}\\
 \int_{t_0}^{t_b}\sup_\theta|R_{11}|\dd t
 &\le C\int_{t_0}^{t_b}|1-\kappa_1|\dd t+o(1),
 &J_G(t)&\le Cq^2,
 \label{an03lcrcv1:rconclusion}\\
 H(t)\le H_{\Gp}(t)&\le\frac98c_b+o(1),
 &E_{\Gp}(t)&\le\frac54c_b+o(1).
 \label{an03lcrcv1:amplitudeconclusion}
\end{align}
All $o(1)$ terms are uniform in time. In particular
$H(t_b)\le(9/8)c_b+o(1)$ is a conclusion. The $\eta=1/8$ clock
has its $o(1)$ defect; $\Gp$ has the bounded defect and relative
transport of RCET. Neither an independent backward clock premise
nor a terminal amplitude cap occurs among the assumptions.
\end{theorem}
\begin{proof}
Use the principal guards
\begin{equation}\label{an03lcrcv1:guards}
 Y<K_Y,\qquad \int_{t_0}^t r<B,\qquad E_{\Gp}<E_g.
\end{equation}
The energy guard is on $E_{\Gp}$, not on $H_{\Gp}$: it bounds
raw weak-coordinate energy before any control of $a/z$ is known.
The cross residual constant $M$ will be chosen from these guards.
For rigorous continuity when the S2 energy is derived rather than
assumed as an instantaneous budget, temporarily stop also at cross
residual sum $Mq$. This auxiliary bound is improved below; it is
neither a retained hypothesis nor a fourth closure route.

\paragraph{Constants chosen before $C_0$.}
Let $C_r,C_x,C_v$ be fixed enlarged constants in RCET Lemma 4.1,
its bound for $r$, and its Volterra inequality. They depend only on
the fixed core and model budgets, not on $C_0,Z$.
Set
\[
 B=C_rB_s+2,\quad E_g=2+2c_b,\quad
 h_*=(9/8)c_b+1,\quad \underline b=(169/400)(1-c_b)>0.
\]
The estimates below will give $E_G\le2H_{\Gp}+Fq^5$, with $F$
fixed only after the geometry is chosen, and the three uniform bounds
\begin{equation}\label{an03lcrcv1:integralconstants}
 \sup E_G\le E_*:=2h_*+1,\quad
 \int E_G\le K_E:=2eh_*/\underline b+1,\quad
 \int\sqrt{E_G}\le K_{1/2}:=
 2\sqrt{2eh_*}/\underline b+1.
\end{equation}
Choose $A_*>1$ larger than
$e^B[1+C_v(2+E_*)K_{1/2}]$ and
$K_Y=4A_*\exp(C_ve^BK_E)$.
Finally choose $M$ more than twice the diagonal norm bound and twice
\[
 C_x\{3+E_g+K_Y\sqrt{E_g+1}\}.
\]
These choices depend on no $C_0$ or geometric cutoff. Now choose
$C_0<C_{\rm adm}(B,M)$ small enough that the conservative
$y$ comparison gives $y\le1/3$ and the $C_0$ part of $D_+$ is
at most $1/2$. Choose RCET's fixed threshold $K=K(C_0)$ and
the resulting finite $Z,Z_\Sigma,R_\Sigma$ and secondary constants.
Finally decrease $q$ to make every fixed geometric constant times
the displayed positive powers of $q$ small. This is the order
$(B,E_g,K_Y,M)$, then $C_0$, then $K,Z$, then $q$.

\paragraph{Estimates on a joint stopped interval.}
Lemma \ref{an03lcrcv1:localclock} gives the sign and defect estimates.
The chart-mass and secondary lemmas give, uniformly up to $t_*$,
\begin{equation}\label{an03lcrcv1:restsmall}
 E_{\rm rest}\le Fq^5+C_\Sigma q^{\Gamma_\nu}H
 \le Fq^5+C_\Sigma q^{\Gamma_\nu}E_g=o(1).
\end{equation}
Here and below $F$ may depend on the now fixed geometric constants.
Also $J_{\rm rest}\le J_G<q^2K_Y^2$ and
$\mu_{\Gp}=E_{\Gp}+J_{\Gp}\le E_g+q^2K_Y^2$.
The pointwise ledger cap therefore implies
\[
 H_{\Gp}(t)\le(9/8)c_b+o(1),\qquad
 E_{\Gp}(t)\le(5/4)c_b+o(1)<E_g-1/2.
\]
In particular $H_{\Gp}(t_*)\le h_*$ follows from the state at
$t_*$ itself, with no reference to $t_b$.

For every $s\le t\le t_*$, the sign-good group's own defect gives
\[
 H_{\Gp}(s)\le eH_{\Gp}(t)e^{-b_0(t-s)},\quad
 \int_{t_0}^{t_*}H_{\Gp}\le eh_*/\underline b,\quad
 \int_{t_0}^{t_*}\sqrt{H_{\Gp}}
 \le2\sqrt{eh_*}/\underline b.
\]
The bound uses bounded distortion, not a vanishing $\eta=0$ defect.
By \eqref{an03lcrcv1:signgood} and \eqref{an03lcrcv1:restsmall},
$E_G\le2H_{\Gp}+Fq^5$ after taking
$C_\Sigma q^{\Gamma_\nu}\le 8/9$.
Since $Fq^5L_T\log(1/q)$ and
$\sqrt Fq^{5/2}L_T\log(1/q)$ tend to zero, all bounds in
\eqref{an03lcrcv1:integralconstants} follow, uniformly in $t_*$.

\paragraph{The unchanged Volterra step and the other guards.}
RCET's instantaneous estimates apply to the whole explicit group $G$:
\begin{align*}
 \frac{\|e_2\|_{P_1}+\|e_1\|_{P_2}}q
 &\le C_x(1+E_G+q^2Y^2+\sqrt{E_G}Y+\mu_{\mathcal B}/q),\\
 r&\le C_r|1-\kappa_1|+
 C_r\{q^2(1+E_G+Y^2+\sqrt{E_G}Y)+\mu_{\mathcal B}\}.
\end{align*}
In particular \eqref{an03lcrcv1:restsmall} gives $E_G\le E_g+1$
under the original guards, so for small $q$ the cross residual sum
is at most $Mq/2$ by the already fixed choice of $M$.
This improves the auxiliary residual stop, without using the
$E_{\Gp}/H_{\Gp}$ comparison to choose $M$.

Reuse RCET Theorem 4.2's Volterra step without changing its kernel:
\begin{equation}\label{an03lcrcv1:volterra}
 Y(t)\le A_*+C_ve^B\int_{t_0}^t E_GY
             +C_ve^Bq^2\int_{t_0}^t\sqrt{E_G}Y^2.
\end{equation}
The choice of $A_*$ dominates the entry term $e^B\sqrt{C_Jq}$
and $C_ve^B\int\sqrt{E_G}(1+E_G+\mu_{\mathcal B}/q)$.
Under $Y<K_Y$, make the last term at most $A_*$ by small $q$.
Gronwall and \eqref{an03lcrcv1:integralconstants} then give
$Y\le2A_*e^{C_ve^BK_E}=K_Y/2$.
Integrating the bound for $r$ gives
\[
 \int_{t_0}^{t_*}r\le C_rB_s+
 O(q^2\log(1/q))+O(\sup\mu_{\mathcal B}\log(1/q))
 =C_rB_s+o(1)<B-1.
\]
Here $\sup\mu_{\mathcal B}=o(q)$ is sufficient.

At entry, $Y=O(q^{1/2})$, $E_{\Gp}=o(1)$ and $\int r=0$.
The chart and secondary entry bounds give $E_{\rm rest}(t_0)=o(1)$,
so RCET's instantaneous estimate also starts the auxiliary residual
bound strictly. All three principal guards and that auxiliary bound
improve with strict slack. Finite-horizon continuity of the norms
and moment integrals, or their standard essential-supremum
continuation form, excludes a first exit before $t_b$.
The displayed estimates thus hold on the entire interval.
Equations \eqref{an03lcrcv1:residualconclusion}--
\eqref{an03lcrcv1:amplitudeconclusion} and the clock claims follow.
\end{proof}

\section{What has been removed, and where this increment stops}
The pointwise response ledger supplies the cap at every stopped
endpoint. Thus the independent backward clock premise in RCET
Theorem 4.2 and its externally supplied terminal amplitude cap are
removed. The backward estimates are now deductions on the same
controlled interval on which they are used.

The result remains conditional on the exact Section 1 budgets.
In particular the original entry and strong-learning estimate are
S1 contracts, and the secondary profile/distance/gain package remains
S1/S2 work. Continued chart geometry is an instantaneous-state
budget; this increment derives its small mass, not its retention.
The constants $Z(C_0)$ in chart cross pressures are harmless:
$\mathfrak F_{\mathcal D}\le C(Z+Z^2q)\mu_{\mathcal D}=o(1)$,
and $E_{\mathcal D}\le Z^2q^2\mu_{\mathcal D}=o(1)$.
More generally, geometric constants from the secondary contracts
multiply only $q^5$ or $q^{\Gamma_\nu}$ here. They are made small
after, not before, fixing the geometry. No smallness of $c_b$ other
than the stipulated fixed inequality $c_b<1$ is required.

All four preliminary checks therefore hold with the stated retained
inputs. Route 3 is complete in this conditional scope. No overlap
profile transfer, $M$-clock passage, capped-gain proof, or bridge
problem is started.

\begin{thebibliography}{9}\small
\bibitem{RCET} RCET, \emph{AN03 residual closure and eta0 transport v1}.
Prefix \texttt{an03rcetv1:}; Theorems 2.1 and 4.2,
Lemma 4.1; labels \texttt{eta0}, \texttt{outputs}, \texttt{closure},
\texttt{volterra}, \texttt{failingendpoint}.
\bibitem{BULK} BULK, \emph{AN03 bulk clock and secondary Q2 energy v1}.
Prefix \texttt{inc:AN03:bulkclock:v1:}; labels
\texttt{bulkclockclosure}, \texttt{moments}, \texttt{gainenergy},
\texttt{secondary}, \texttt{secondarydetail}, \texttt{gammamargin},
\texttt{transverse}, \texttt{residualinputs}.
\bibitem{CLOCK} CLOCK, \emph{AN03 exact amplitude clock and c passage v1}.
Prefix \texttt{inc:AN03:clock:v1:}; labels \texttt{identity},
\texttt{endpoint}, \texttt{endpointidentity}, \texttt{Hcap}.
\bibitem{F} \emph{RELU population foundations v2}, P5:
balance, characteristic regularity and loss dissipation.
\bibitem{ROADMAP} \emph{Unified analytic theory roadmap}, v1.1,
October 6, 2026: S0 Route 3 and ledger items 9--11.
\end{thebibliography}
\end{document}
